% ════════════════════════════════════════════
% 摘要（00.摘要.tex）—— 写作提示词见本文件末尾【本章节写作提示词】。
% ════════════════════════════════════════════
\begin{abstract}

本文针对低空空域小型无人机快速测向定位问题，结合空间几何约束判定、交会误差传播、组合调度与动态重调度等方法，提出了多移动观测车协同定位与调度的系统解决方案。

\textbf{针对问题一:} 本问是单目标可观测区域判定与单目标连续优化问题，需在车辆几何、视线遮挡与云台俯仰三项约束下确定目标 T1 的可观测区域并求最优观测站。首先，将车辆等效为半径 1.5 m 的圆盘，建立区域边界与建筑不侵入约束；接着，用三维线段与长方体棱柱的 slab 相交判定视线遮挡；然后，以俯仰角范围约束建立水平距离下界，并以 Q 等于 r 的倒数定义观测质量，将最优站问题化为最小水平距离问题；最后，解析与 0.5 m 网格数值联合求解。结果表明，可观测区域面积约 5697 平方米，最优水平距离为 7.582 m，俯仰角恰为 80 度，最优站为 91.155 米与 66.121 米处，就近车辆 V4 行驶 3.67 m 即可到达。

\textbf{针对问题二:} 本问是两站交会定位精度建模与多目标车辆调度优化问题，需确定告警 T2 的最优站对并完成三告警并发下的车辆复用与路径规划。首先，由方位角测量误差服从零均值、标准差 1 度的正态分布，线性化推导位置误差均方根等于标准差乘两站距离平方和的平方根再除以交会角正弦；接着，以精度指标为正弦交会角除以上述平方根，在行驶不超过 30 s 约束下枚举可观测站对；然后，以各告警可观测站与交会角约束建立车辆指派与时序仿真模型，最小化最大观测开始时间；最后，与独立定义的最近指派启发式对比。结果表明，T2 最优站对为西站与北站，基线 8.229 m，交会角 90 度，定位误差 14.36 cm；并发调度最优最大观测开始时间为 12.063 s，优于启发式的 12.561 s。

\textbf{针对问题三:} 本问是定位误差与指向误差传播的概率约束问题，需确定满足方位偏差不超过 5 度概率不小于 0.95 的最大行驶距离及误差模型参数约束。首先，由车辆定位误差与云台零位偏差、控制误差的独立性建立方位偏差传播模型；接着，线性化得到方位偏差标准差等于定位标准差除以水平距离；然后，按 95\% 分位数反解定位标准差上限；最后，以 20 万次蒙特卡洛验证小角度近似。结果表明，最大允许行驶距离为 66.92 m，行驶 80 m 时需满足 a 加 80b 不大于 1.638 m，定位误差、零位偏差、控制误差的方差贡献约为 84.6\%、9.8\% 与 5.5\%。

\textbf{针对问题四:} 本问是扰动下的动态重调度与方案比较问题，需评估新增告警、随机车辆故障与通信中断及组合扰动下的重调度收益。首先，以问题二所得调度方案为基线建立 t 等于 5 s 时刻的车辆状态快照；接着，建立重调度枚举模型与保持原方案策略，在 60 s 评价窗口内按完成率优先、最晚观测开始时间次优排序；然后，设计 8 个新增告警方向、5 个故障车与 3 个通信恢复时刻及 4 个组合场景共 20 个代表性场景；最后，逐场景比较。结果表明，重调度在故障与组合扰动下可将完成率由 0.67 或 0.75 提升至 1.00。

综上，本文完成了可观测区判定、交会定位、并发调度、误差传播与扰动重调度全流程研究。研究成果可为低空空域多车协同探测与指挥调度提供理论依据和决策参考。

\keywords{无人机定位；多站交会；误差传播；车辆调度；动态重调度}

\end{abstract}
